%\frontmatter
%\newpage
\thispagestyle{empty}
\vspace*{4cm}
\begin{center}
\Huge\textsc{Tattvayogabindu}
\end{center}
\vfill

\newpage %%%%%%%%%%%%%%%%%%%%%%%%%%%%%%%%%%%
\thispagestyle{empty}

\scriptsize
\noindent
\textsl{L'Institut Français de Pondichéry (IFP), UMIFRE 21 CNRS-MAE, est un établissement à autonomie financière sous la double tutelle du Ministère français des Affaires Etrangères (MAE) et du Centre National de la Recherche Scientifique (CNRS). Il est partie intégrante du réseau des 27 centres de recherche de ce Ministère. Avec le Centre de Sciences Humaines (CSH) à New Delhi, il forme l'USR 3330 du CNRS «~Savoirs et Mondes Indiens~». Il remplit des missions de recherche, d'expertise et de formation en Sciences Humaines et Sociales et en Ecologie en Asie du Sud. Il s'intéresse particulièrement aux savoirs et patrimoines culturels indiens (langue et littérature sanskrite, histoire des religions, études tamoules…), aux dynamiques sociales contemporaines, et aux écosystèmes naturels de l'Inde du Sud.}

\medskip
\noindent
The French Institute of Pondicherry (IFP), UMIFRE 21 CNRS-MAE, is a financially autonomous institution under the joint supervision of the French Ministry of Foreign Affairs (MAE) and the French National Centre of Scientific Research (CNRS). It is a part of the network of 27 research centres under this Ministry. It also forms part of the research unit 3330 “Savoirs et Mondes Indiens” of the CNRS, along with the Centre de Sciences Humaines (CSH) in New Delhi. It fulfils its missions of research, expertise and training in Human and Social Sciences and Ecology in South Asia. It works particularly in the fields of Indian cultural knowledge and heritage (Sanskrit language and literature, history of religions, Tamil studies…), contemporary social dynamics and the natural ecosystems of South India.

\medskip\noindent
French Institute of Pondicherry, \\
11, St.\,Louis Street, P.B.\,33, \\
Pondicherry---605001, India.\\
Tel: +91 413 2231600 \\
Email: ifpdir@ifpindia.org \\
Website: \url{https://www.ifpindia.org}

\vfill
\begin{center}
\includegraphics[scale=.16]{flourish.jpg}
\end{center}
\vfill

\noindent
\textsl{L'École française d'Extrême-Orient (EFEO), fondée en 1900 à Hanoï, est un établissement relevant du ministère français de l'Enseignement supérieur et de la Recherche dont la mission scientifique est l'étude des civilisations classiques de l'Asie. Son champ de recherches s'étend de l'Inde à la Chine et au Japon et, englobant l'ensemble du Sud-Est asiatique, comprend la plupart des sociétés qui furent indianisées ou sinisées au cours de l'histoire. Autour de ses dix-sept centres et antennes, installés dans douze pays d'Asie, se sont constitués des réseaux de chercheurs locaux et internationaux sur lesquels l'École a pu s'appuyer pour construire son essor. L'EFEO aborde l'Asie par des recherches pluridisciplinaires et comparatistes, associant l'archéologie, l'histoire, l'anthropologie, la philologie, et les sciences religieuses. À Pondichéry, les projets de l'EFEO portent essentiellement sur l'«~indologie~» classique~: sanskrit, tamoul ancien, histoire, histoire de l'art et des religions.}

\medskip
\noindent
The mission of The French School of Asian Studies (EFEO), founded in 1900 in Hanoi and today under the aegis of the French Ministry of Higher Education and Research, is to study the classical civilisations of Asia. Stretching from India, in the West, across the whole of South-East Asia to China and Japan, the EFEO's research areas cover most of the societies which have been `Indianised' or `Sinicised' over the course of history. A network of international scholars working at the EFEO's seventeen centres and branch offices, which are spread across twelve Asian countries, has been essential in the development of the School's research programme. Interdisciplinary projects bring together leading scholars in the fields of anthropology, archaeology, history, philology, and religious studies. In Pondicherry, the projects of the EFEO focus mainly on classical Indology: Sanskrit, Old Tamil, History, and History of art and of religions.

\medskip
\noindent
École française d'Extrême-Orient,\hfill Pondicherry Centre of the EFEO\\
22, avenue du Président Wilson,  \hfill 16 \&\ 19, Dumas Street, \\
75116 Paris, France.             \hfill Pondicherry---605 001, India. \\
Tel: +33 1 53 70 18 60          \hfill Tel: +91 413 2334539/2332504\\
Website: \url{https://www.efeo.fr}    \hfill Email: administration@efeo-pondicherry.org


\newpage %%%%%%%%%%%%%%%%%%%%%%%%%%%%%%%
\thispagestyle{empty}
\small
\begin{center}
\textbf{THE HAṬHA YOGA SERIES}
\end{center}

Yoga is central to Indian religious practice and culture. From probable origins among heterodox ascetics in the first millennium \textsc{bce} it gradually became part of almost all of India’s religious traditions. Key to yoga’s importance in both its modern globalised manifestations and Indian religious practice are its physical techniques. Some ascetic physical practices are as ancient as yoga itself, while others appear to be innovations introduced at the beginning of the second millennium \textsc{ce} when a corpus of works on \textit{haṭha} — as the method of yoga in which physical practices predominate is known in Sanskrit — was composed. 	

The \textit{Haṭha Yoga Project} based at SOAS University of London (2015–2020, European Research Council Grant n°\,647963), sought to improve the textual foundations for the study of Haṭha Yoga by critically editing ten of its most important texts, both from the period of formalisation, but also from key moments in its subsequent development up to the 19th century, and to supplement textual evidence with ethnographic observation of its ascetic practitioners in India today. This mini-series, launched within the ‘Collection Indologie’, will publish some of the fruits of the project, in particular editions and translations of Sanskrit manuals of Haṭha Yoga.

The \textit{Haṭha Yoga Series} is not a closed collection: as well as the list of works of which 
editions and translations were promised as part of the project, studies of numerous related works were inspired or further advanced during the project's workshops, and some of these will also be published in the series.
%\normalsize

\cleardoublepage %%%%%%%%%%%%%%%%%%%%%%%%%%%%%%%%%%%
\thispagestyle{empty}
\begin{center}
\scshape
Collection Indologie ???\\
Haṭha Yoga Series ?

\vskip 1in

\Huge 
The Tattvayogabindu of~Rāmacandra

\vskip .5in

\LARGE 
A Critical Edition and Annotated~Translation of an Early Modern Text on Rājayoga, with a Comparative Analysis of the Complex Yoga Taxonomies from the Same Period\\

\vskip .9in

\Large 
Nils Jacob Liersch

\vfill

\normalsize
institut français de pondichéry\\
école française d’extrême-orient
\end{center}
\upshape

%%%%%%%%%%%%%%%%%%%%%%%%%%%%
\newpage
\thispagestyle{empty}
%\begin{absolutelynopagebreak}
{\footnotesize
\noindent
Comité de rédaction / Editorial Board

\smallskip\noindent
Hugo \textsc{David},
T.~\textsc{Ganesan},
Dominic \textsc{Goodall},
Borayin \textsc{Larios},
Blandine \textsc{Ripert}

\bigskip\noindent
Comité Éditorial / Advisory Board 
 
\smallskip\noindent
Diwakar \textsc{Acharya} (All Souls College, Oxford University),\\ 
Nalini \textsc{Balbir} (Université de Paris III et École pratique des hautes études),\\ 
Peter \textsc{Bisschop} (Leiden University),\\ 
Eloïse \textsc{Brac de la Perrière} (Université de Paris IV),\\
Sylvain \textsc{Brocquet} (Université d’Aix-Marseille),\\
Whitney \textsc{Cox} (Chicago University),\\
Richard \textsc{Davis} (Bard College, New York),\\
Arlo \textsc{Griffiths} (École française d’Extrême-Orient),\\ 
Oskar \textsc{von Hinüber} (University of Freiburg im Breisgau),\\ 
Padma \textsc{Kaimal} (Colgate University),\\ 
Kei \textsc{Kataoka} (Kyushu University),\\ 
Vempati \textsc{Kutumba Sastry} (Banaras Hindu University),\\ 
Leslie \textsc{Orr} (Concordia University),\\ 
Parul \textsc{Pandya Dhar} (Delhi University),\\
Aloka \textsc{Parasher-Sen} (University of Hyderabad),\\ 
V.\ \textsc{Selvakumar} (Thanjavur University),\\
Kesavan \textsc{Veluthat} (Delhi University, retired).

\bigskip\noindent
Comité de Lecture / Peer-review

\smallskip\noindent
Les éditeurs font appel à des spécialistes de leur choix pour l’évaluation des manuscrits soumis.\\
The publishers of the series call on experts of their choice for the evaluation of manuscripts submitted.

\vfill

\noindent
\copyright\ Institut Français de Pondichéry, 2026 (ISBN 978-81-8470-???-?)\\
\copyright\ École française d’Extrême-Orient, 2026 (ISBN 978-2-85539-???-?)\\

\bigskip\noindent
Typeset by Nils Jacob \textsc{Liersch}\\
Cover design: Sri Aurobindo Press, Pondicherry.\\
Cover photo: ??\\
Printed at the Sri Aurobindo Ashram Press, Pondicherry.}
%\end{absolutelynopagebreak}
